\documentclass[12pt, letterpaper]{article}
\usepackage[margin=1in]{geometry}
\usepackage{amsmath, amssymb, amsthm}
\usepackage{graphicx}
\usepackage{hyperref}
\usepackage{authblk}

\hypersetup{
    colorlinks=true,
    linkcolor=blue,
    filecolor=magenta,      
    urlcolor=cyan,
    pdftitle={An Algebraic System Based on the Golden Ratio and its Application to Complex Series},
    pdfauthor={Tristen Harr},
}

% Theorem-like environments
\newtheorem{theorem}{Theorem}[section]
\newtheorem{lemma}[theorem]{Lemma}
\newtheorem{definition}[theorem]{Definition}
\newtheorem{conjecture}[theorem]{Conjecture}
\newtheorem{corollary}[theorem]{Corollary}

% Custom command for the main operator
\newcommand{\LambdaG}{\Lambda_{G1}}

\title{\textbf{An Algebraic System Based on the Golden Ratio and its Application to Complex Series}}
\author{Tristen Harr}
\affil{\textit{Saint Charles, Missouri}}
\date{June 18, 2025}

\begin{document}
\maketitle

\begin{abstract}
This paper defines and explores a novel algebraic system, termed the "Golden Algebra," whose structure is axiomatically determined by the golden ratio $\phi$. We derive the algebra's core constants and operators, and prove its fundamental identities. The richness of this system is demonstrated by deriving closed-form solutions for several families of infinite series involving the algebra's primary operator, $\Lambda_{G1}$. These results include master theorems for series weighted by polynomial and reciprocal powers, the latter showing a native connection to the Polylogarithm special function. We also prove an identity connecting an infinite series of Zeta function values to the cotangent function. Finally, we observe a structural parallel between a potential function arising from this algebra and the functional equation of the completed Riemann Zeta function, $\xi(s)$. This observation leads us to propose a new conjecture: that the non-trivial zeros of $\zeta(s)$ behave as stable resonances within this algebraic framework, a conjecture which would imply the Riemann Hypothesis.
\end{abstract}

\section{Axiomatic Foundation of the Golden Algebra}
\subsection{Axiomatic Definition of the Constants}
We define a complex number field, the "Golden Algebra," based on the components ($a, b$) of a generic 2D rotation-scaling matrix $\begin{pmatrix} a & -b \\ b & a \end{pmatrix}$. The system is defined by two fundamental axioms:
\begin{enumerate}
    \item \textbf{Axiom 1 (Additive):} $a+b=1/2$.
    \item \textbf{Axiom 2 (Ratio):} $a/b=\phi$, where $\phi$ is the Golden Ratio.
\end{enumerate}
This system of equations yields a unique solution for the real constants $a$ and $b$: $a=(\sqrt{5}-1)/4$ and $b=(3-\sqrt{5})/4$. We designate these as the algebra's primary constants, T and J, respectively. All subsequent properties of the algebra are consequences of these two axioms.

\subsection{Core Identities and Operators}
\begin{theorem}[Core Algebraic Identities]
The constants T and J satisfy the following relations:
\begin{enumerate}
    \item $T+J=1/2$ 
    \item $T/J=\phi$ 
\end{enumerate}
\end{theorem}
\begin{proof}
The identities are a direct consequence of the axioms, proven by substituting the derived values for T and J. For example, $T+J=\frac{\sqrt{5}-1}{4}+\frac{3-\sqrt{5}}{4}=\frac{2}{4}=1/2$.
\end{proof}

From these constants, we define the primary operator of the algebra, $\LambdaG = T+iJ$.

\begin{lemma}
The operator $\LambdaG$ is convergent.
\end{lemma}
\begin{proof}
Convergence requires $|\LambdaG|^2 < 1$. We calculate $|\LambdaG|^2 = T^2+J^2 = \left(\frac{\sqrt{5}-1}{4}\right)^2 + \left(\frac{3-\sqrt{5}}{4}\right)^2 = \frac{20-8\sqrt{5}}{16} = \frac{5-2\sqrt{5}}{4} \approx 0.13 < 1$. The condition is satisfied.
\end{proof}

\section{Analysis of Series within the Golden Algebra}
We now demonstrate the utility of this algebraic system by analyzing several families of infinite series involving the operator $\LambdaG$.

\subsection{Theorem for Polynomial-Weighted Series}
\begin{theorem}[Polynomial-Weighted Sums]
The infinite series $S_k = \sum_{n=0}^{\infty} n^k \cdot (\LambdaG)^n$ for any integer $k \ge 0$ converges to the exact value: $$S_k = \frac{A_k(\LambdaG)}{(1-\LambdaG)^{k+1}}$$ where $A_k(x)$ is the k-th Eulerian polynomial, whose coefficients can be generated by the recurrence relation $A_{k+1}(x) = x(1-x)A'_k(x) + x(k+1)A_k(x)$ with base case $A_0(x)=1$.
\end{theorem}
\begin{proof}
This is a known result for geometric series. The proof follows by recursively applying the operator $x \frac{d}{dx}$ to the geometric series $S_0(x) = \frac{1}{1-x}$. The convergence of our operator $\LambdaG$ ensures the validity of this method.
\end{proof}

\subsection{Theorem for Reciprocal-Weighted Series}
\begin{theorem}[Connection to the Polylogarithm Function]
The series $S_s = \sum_{n=1}^{\infty} \frac{(\LambdaG)^{n}}{n^s}$ converges for any integer $s \ge 1$ to the exact value: $$S_s = \mathrm{Li}_s(\LambdaG)$$ where $\mathrm{Li}_s(z)$ is the Polylogarithm function of order s.
\end{theorem}
\begin{proof}
The proof is by mathematical induction on the order $s$. The base case for $s=1$ is the series for the complex logarithm, proven by integrating the geometric series. The inductive step involves dividing the identity for order $s$ by $x$ and integrating from $0$ to $t$, which is a standard method that directly yields the identity for order $s+1$.
\end{proof}

\subsection{A Connection to Analytic Number Theory}
\begin{theorem}[The Harmonic Cotangent Identity]
The following identity, relating an infinite series of Zeta function values to $\pi$ and the cotangent function, holds for any integer $n\ge1$:
$$\sum_{k=1}^{\infty}\frac{(n^{k}-1)\zeta(k+1)}{(n+1)^{k}}=\pi \cot\left(\frac{\pi}{n+1}\right)$$
\end{theorem}
\begin{proof}
This identity is a consequence of a known identity for the Polygamma function: $\sum_{k=1}^{\infty} \frac{(n^k-1)\zeta(k+1)}{(n+1)^k} = \psi^{(0)}(\frac{n}{n+1}) - \psi^{(0)}(\frac{1}{n+1})$. The theorem is then proven by applying the standard reflection formula for the Polygamma function, $\psi^{(0)}(1-z)-\psi^{(0)}(z)=\pi \cot(\pi z)$, with the substitution $z = \frac{1}{n+1}$.
\end{proof}

\section{A Conjectured Correspondence with the Riemann Zeta Zeros}
The rich structure of the Golden Algebra motivates an inquiry into its potential relevance to unsolved problems. Here we propose a speculative connection to the non-trivial zeros of the Riemann Zeta function.

\subsection{A Stability Potential from the Algebra}
\begin{definition}[Symmetric Potential]
We define a potential function $V(s)$ for a system balanced between a parameter $s$ and its reflection $1-s$ using the algebra's constants: $V(s)=T+sJ+(1-s)K$, where $K:=-1/2-T$.
\end{definition}

\begin{theorem}
The potential function $V(s)$ simplifies to the identity $V(s)=s-1/2$.
\end{theorem}
\begin{proof}
The potential can be rearranged as $V(s)=(T+K)+s(J-K)$. Substituting the definitions of the constants, we have $T+K=T+(-1/2-T)=-1/2$ and $J-K=(1/2-T)-(-1/2-T)=1$. Thus, $V(s)=-1/2+s(1)=s-1/2$.
\end{proof}

\subsection{The Central Conjecture}
The completed Zeta function $\xi(s)$ is known to obey the functional equation $\xi(s) = \xi(1-s)$, which implies its derivative has the reflectional anti-symmetry $\xi'(s)=-\xi'(1-s)$. This structural feature is shared by our potential function, $V(s)=s-1/2$. This shared symmetry motivates our central conjecture.

\begin{conjecture}[The Harmonic Resonance Conjecture]
The non-trivial zeros of the Riemann Zeta function, $\rho$, behave as stable resonances in a system governed by the Golden Algebra's stability potential, $V(s)$. The condition for such a stable resonance is that it must lie where the real part of the potential is zero: $Re(V(\rho))=0$.
\end{conjecture}

An immediate consequence of this conjecture, if it were true, is the Riemann Hypothesis. The condition $Re(V(\rho))=0$ implies $Re(\rho - 1/2) = 0$, which requires $Re(\rho) = 1/2$.

\begin{figure}[htbp]
    \centering
    \includegraphics[width=0.8\textwidth]{energy_landscape.png}
    \caption{The potential landscape defined by $V(s)$. The plot illustrates the absolute value of the potential, $|\text{V}(\sigma)|=|\sigma-1/2|$, for real $\sigma$. The potential has a unique minimum at $\sigma=1/2$. Our conjecture proposes that the Zeta zeros must lie in this valley.}
    \label{fig:energy_landscape}
\end{figure}

\subsection{Evidence in Support of the Conjecture}
While not a proof, we can test the conjecture for consistency against known data. The locations of the trillions of calculated non-trivial zeros, $\{\rho_{n}\}$, are all of the form $\rho_{n} = 1/2 + it_n$. For every known zero, the condition $Re(V(\rho_n)) = Re(1/2 + it_n - 1/2) = 0$ is satisfied. This perfect agreement between the known data and the conjecture provides strong motivation for its further study.

\section{Conclusion}
We have defined and analyzed a mathematical system, the "Golden Algebra," whose properties are derived from axiomatic constraints involving the golden ratio. We have shown this system is mathematically rich, providing a framework for deriving closed-form solutions for several families of infinite series and connecting to established special functions.

Based on an observed structural symmetry, we have proposed a new, falsifiable conjecture: that the non-trivial zeros of the Riemann Zeta function are governed by the stability principles of this algebra. The primary contribution of this paper is the formalization of the Golden Algebra and the clear statement of this conjecture. Proving or disproving this conjectured link remains an open problem for future work, but the consistency with all known data suggests it may be a fruitful avenue of inquiry.

\end{document}